\section*{Background and Significance}

A selective classifier issues a decision only when its confidence clears a
cutoff and refers the remaining cases to a human reader. For chest
radiographs, the per-finding decision thresholds and the abstention cutoff are
normally chosen on data from the developing hospital. If the system is deployed
elsewhere unchanged, the operational question is whether that frozen policy
still delivers the error rate and referral rate it was calibrated for.
Re-fitting thresholds at the new site answers a different question: whether the
method can be re-tuned.

Multimodal models add a second shift. Models that read the clinical indication
or history alongside the image depend on text whose availability and content
reflect local documentation. Institutional and contextual shift therefore
travel together, and an evaluation that varies only one cannot estimate how
they combine.

Each component has been studied separately. Selective classification has been
extended to distribution shift on non-clinical benchmarks.\cite{liang2024}
Failure-detection evaluation practice has been examined, finding simple
confidence baselines hard to beat and proposing the area under the generalised
risk--coverage curve (AUGRC).\cite{jaeger2023,traub2024,bernhardt2022}
Rejection mechanisms have been added to multi-label chest-radiograph
classifiers evaluated within datasets.\cite{aperstein2025} Clinical indication
and history improve chest-radiograph classification within a single
source,\cite{jacenkow2022,yang2025} and missing modalities have been treated as
an architectural problem.\cite{wang2025} Cross-dataset failure of
chest-radiograph models has been attributed to shift in labels more than in
images,\cite{cohen2020} and report-derived labels are imperfect proxies for
image findings.\cite{jain2021} Using the same labeller, such as
CheXbert,\cite{smit2020} at two sites does not guarantee equal labelling
accuracy or equal label availability.

We found no study that carries a frozen selective policy to a second hospital
while manipulating pre-diagnostic context under preregistered rules. We also
found none that isolates, with predictions and thresholds held fixed, how label
availability and handling determine the transported risk. This paper does
both. It reports a preregistered compound-shift evaluation, then an audit,
carried out after the external results were seen, which found that the
originally reported cross-site degradation was an artifact of the selective-risk
estimator. The significance is methodological and practical. Cross-site
selective-risk benchmarks built on report-derived labels can reach opposite
conclusions under defensible analysis choices, and we show how to detect and
bound this.
